\section*{T\_QG\_NoGo: A Bounded-Grammar No-Go Theorem for the Fused/Derived QG Object}

\paragraph{Grade.}
This is a finite-carrier bounded-grammar theorem.  It consolidates Lemma 1
(Step 30) and Lemma 2 (Step 31).  It is not an unconditional analytic theorem
and not an absolute statement about every possible physical theory.

\subsection*{Definitions}

\paragraph{Bounded grammar \(G^\ast\).}
In this theorem, \(G^\ast\) is the finite grammar consisting of:
\begin{itemize}
\item the native carrier \(L=(d0,d1,d2,d3)\);
\item endpoint access quotients
  \(q_{QM}=(d0,d1,d2)\) and \(q_{GR}=(d0,d2,d3)\);
\item the finite field-readout model in which a field \(\psi\) supplies
  \(\mathrm{Born}[\psi]=(\rho,j)\) and \(T[\psi]=(T00,T0i)\);
\item the native promotion-bridge gates
  \(G_{suff},G_{desc},G_{stab},G_{ctrl},G_{nosmuggle},G_{vis},G_{audit},G_{strict}\);
\item the anti-reductionism guardrail: a lawful move may chart, type,
  recognize, construct, simulate, or audit a layer, but may not smuggle one
  endpoint as a derivation of the other across a foreclosed boundary.
\end{itemize}

\paragraph{Fused/derived QG object.}
The forbidden target tested here is not the co-sourcing common-refinement
package.  It is a single package that is intended to be both endpoint roles at
once by one of two routes:
\begin{enumerate}
\item \emph{directed route}: one endpoint is a lawful quotient or derivation of
  the other;
\item \emph{fused route}: a single object is formed by stapling the endpoint
  packages together or by one audit serving both endpoint roles.
\end{enumerate}

\paragraph{Lawful in \(G^\ast\).}
A candidate is lawful in \(G^\ast\) only if its quotient maps, audits, route
orders, and controls pass the declared native gates and no-smuggling checks on
the finite carrier.

\subsection*{Theorem}

Within \(G^\ast\), no single lawful fused/derived endpoint object forms by the
directed or fused routes.

More precisely:
\begin{enumerate}
\item The directed route fails: neither endpoint quotient factors through the
  other, endpoint audits are not mutually recovered by the finite held-out
  diagnostic, and the two directed route completions do not commute.
\item The fused route fails: the direct-sum object fails \(G_{nosmuggle}\), and
  the field-derived Born and stress-energy audits do not agree as one shared
  audit on the density/transport overlap.
\end{enumerate}

The co-sourcing common-refinement package constructed in Steps 18--28 remains
the lawful reconciliation candidate in this finite grammar.  Full uniqueness
of that reconciliation is a separate Part-B obligation and is not claimed by
this theorem.

\subsection*{Proof}

\paragraph{Directed route.}
Step 30 computes the quotient factorization defects:
\[
|\Delta(q_{GR}\mid q_{QM})|=8,
\qquad
|\Delta(q_{QM}\mid q_{GR})|=8.
\]
Thus neither endpoint quotient factors through the other on the declared
carrier.

Step 30 also computes audit held-out residuals:
\[
\mathrm{res}(T[\psi]\leftarrow(\rho,j))=0.5991748529352457,
\]
and
\[
\mathrm{res}((\rho,j)\leftarrow T[\psi])=0.5776762928080916.
\]
The directed route-order diagnostic gives normalized mismatch \(0.5\).

The vertical-stack control in Step 30 has factorization defect \(0\) and route
mismatch \(0.0\).  Therefore the test is not rigged to reject every directed
stack.  Source:
\texttt{steps/step30\_qg\_directed\_reduction\_nogo\_artifacts/step30\_results\_summary.md}.

\paragraph{Fused route: union.}
Step 31 computes the direct-sum map
\[
\pi_{union}=(d0,d1,d2,d0,d2,d3).
\]
It has target dimension \(6\), rank \(4\), extra coordinate count \(2\), and
duplicate pair count \(2\).  Thus it fails \(G_{nosmuggle}\).  The
nonduplicating refinement control has dimension \(4\), rank \(4\), extra
coordinate count \(0\), and passes.  Source:
\texttt{steps/step31\_qg\_fused\_object\_nogo\_artifacts/step31\_results\_summary.md}.

\paragraph{Fused route: shared audit.}
Step 31 uses the Step 24 field-derived audits.  On the shared
density/transport overlap, the joint equal residual is
\[
0.9935981568901483.
\]
The density proportional residual is \(0.5207123913178113\), and the
transport proportional residual is \(0.9994686548331937\).  Hence the single
shared-audit route fails on the declared finite data.  The agreeing-audit
control has residual \(0.0\) and passes.  Sources:
\texttt{steps/step24\_derived\_physical\_audits\_artifacts/step24\_results\_summary.md}
and
\texttt{steps/step31\_qg\_fused\_object\_nogo\_artifacts/step31\_results\_summary.md}.

\paragraph{Constructed lawful reconciliation candidate.}
Steps 18--22 classify \(L\) as an admissible joint quotient, a computed
\(\texttt{BridgeMediatedRole}\), scale-consistent on the toy, and an F51
common-refinement unification on the finite carrier.  Steps 24--28 add the
finite physical content: \(\psi\) co-sources \(\mathrm{Born}[\psi]\) and
\(T[\psi]\), admits a self-consistent semiclassical toy dynamic, and admits
stationary plus non-stationary relational clock-field representations.
Sources:
\texttt{steps/step18\_f37\_complementarity\_f24\_resolution\_artifacts/step18\_results\_summary.md},
\texttt{steps/step22\_f51\_unification\_common\_refinement\_artifacts/step22\_results\_summary.md},
and
\texttt{steps/step29\_consolidated\_statement\_artifacts/consolidated\_findings\_step29.md}.

\paragraph{Conclusion.}
The directed and fused routes fail, each with a passing can-fail control.
Therefore the fused/derived endpoint object is not lawful in \(G^\ast\).  The
standing lawful candidate is the co-sourcing common-refinement, with uniqueness
left to Part B.

\subsection*{Seven-Gate Audit}

\begin{enumerate}
\item Self-contained definitions: \(G^\ast\), directed route, fused route, and
  lawful candidate are defined above.
\item Route evidence: directed, union, and shared-audit routes cite Step 30 and
  Step 31 computations.
\item Controls: vertical-stack, nonduplicating-refinement, and agreeing-audit
  controls pass.
\item Bounded grade: finite-carrier bounded-grammar theorem.
\item No-smuggling: \(G_{nosmuggle}\) is explicitly applied to the union route.
\item Scope boundary: no statement beyond \(G^\ast\) is asserted.
\item Next obligation: Part B must test uniqueness or exhaust lawful
  reconciliation families.
\end{enumerate}

\subsection*{Interpretation}

The framework-side interpretation is that the technical walls are the grammar
enforcing its no-smuggling and anti-reduction commitments.  P3 appears as
route-order mismatch; P2 appears as the requirement for constrained relational
time rather than external-time presentation.  The lawful object recovered in
the finite construction is not a fused endpoint object; it is the co-sourcing
field layer.  See
\texttt{interpretation/quantum\_gravity\_is\_illegal.md}
and
\texttt{interpretation/ladder\_vs\_fork.md}.

\subsection*{Nonclaim}

This theorem is a bounded statement internal to \(G^\ast\).  It does not make
an absolute claim about all possible physical theories, does not assert a
complete physical model, and does not close the Part-B uniqueness program.
